% ============================================================================
%  chapter1.tex
%  The Foundations of Memory and Pointers
% ============================================================================

\chapter{The Foundations of Memory and Pointers}

\section{The Machine Underneath}

Before we talk about C, we must talk about the machine. Everything a
computer does reduces to a staggeringly simple model: there is a large
array of bytes, each byte has a number (its \emph{address}), and the
processor runs a loop that reads instructions from memory and acts on
them. The CPU has no understanding of ``variables'', ``strings'', or
``objects''. It understands \emph{addresses} and \emph{bytes}. When we
write C, we are simply describing, with a thin layer of sugar, what the
machine will actually do.

This is the heart of the Sanfilippo approach to C: we never hide the
machine from the student. If you understand bytes, addresses, and
lifetimes, then every higher-level concept---data structures, network
servers, databases, virtual machines---is just a disciplined arrangement
of these three ideas. If you skip this foundation, every later chapter
will feel like magic, and magic breaks in production.

Let us begin with the smallest unit of storage the CPU can address
directly: the \emph{byte}.

\begin{verbatim}
   A byte is 8 bits. Each bit is a switch: 0 (off) or 1 (on).

   Address 0x7ffe29a00000:

        +---+---+---+---+---+---+---+---+
        | 0 | 1 | 0 | 0 | 0 | 0 | 1 | 1 |   = 0x43 = 67 = ASCII 'C'
        +---+---+---+---+---+---+---+---+
         bit7 bit6 bit5 bit4 bit3 bit2 bit1 bit0

        bit0 has value 1 (2^0), bit1 has value 2 (2^1),
        bit2 has value 4 (2^2), and so on, up to bit7 = 128.
        This byte represents the number 64 + 2 + 1 = 67.
\end{verbatim}

Every byte in the address space has a unique address. On x86-64 the
implemented virtual address space is actually numbered from $0$ to
$2^{48}-1$ (or to $2^{57}-1$ with five-level paging); the 64-bit register width
would allow $2^{64}$, but the hardware does not implement that full range. Either
way, it is so large that we will never fill it. We write addresses in hexadecimal
because it maps cleanly onto groups of four bits, and because writing
$2^{64}$ as decimal numbers is a waste of human attention.

\begin{insight}[Why hexadecimal]
A byte is two hexadecimal digits. One hex digit is four bits---a
\emph{nibble}, half a byte. That is
why \code{0x43} reads like a byte: \code{0x} is the prefix, \code{4} is
the high nibble, \code{3} is the low nibble. When you see a memory
dump, you can mentally group it into bytes by pairs of digits.
\end{insight}

\subsection{Bits, Bytes, and Words}

Three sizes will keep recurring, so let us settle them now. A
\emph{bit} is one switch, 0 or 1. A \emph{byte} is eight bits and is
the smallest addressable unit. A \emph{word} is the machine's
\emph{natural unit of data}: the size of a general-purpose register,
and the size the CPU prefers to load and store in one operation.

\begin{verbatim}
   The three units, drawn to scale:

     a bit:        a byte (8 bits):        a word (a register):
                   1 0 1 1 0 0 1 0
     +---+          +---+---+---+---+---+---+---+---+       +---------------+
     | 1 |          | 1 | 0 | 1 | 1 | 0 | 0 | 1 | 0 |       |   0x00000043  |
     +---+          +---+---+---+---+---+---+---+---+       +---------------+
                                            ^
                                            ^   four bytes grouped:
                                                the natural 64-bit word
                                                = 0x0000000000000043

   On a 64-bit machine, one word is 8 bytes and holds up to 2^64
   different values -- enough to address every byte in the address
   space. That is no coincidence: a POINTER is exactly one word.
\end{verbatim}

The word size is the machine's identity, and it has consequences you
will feel all book long:

\begin{itemize}
  \item A \emph{pointer} is one word. On a 32-bit machine pointers are
    4 bytes; on a 64-bit machine they are 8. That is why
    \code{sizeof(int *)} printed 8 in Chapter~1's table, and why the
    address space is $2^{64}$ bytes on modern systems.
  \item The CPU likes \emph{aligned} access: reading a word from an
    address divisible by the word size is fastest; straddling a
    boundary costs extra. Alignment, and the padding it forces into
    structs, is the subject of Chapter~3.
  \item The word survives in our vocabulary: a \emph{double word} is
    two words, hex dumps are laid out in words, and ``word size'' is
    how we say ``how wide is this machine?''.
\end{itemize}

\begin{verbatim}
   Quick sanity check -- every claim above in one line:

     64-bit machine   word = 8 bytes   pointer = 8 bytes
     32-bit machine   word = 4 bytes   pointer = 4 bytes
     sizeof(int *) always equals sizeof(void *) on any one machine:
     both are one word, whatever that word is.
\end{verbatim}

\begin{insight}[Why we say "byte" and not "word" for addresses]
The CPU addresses memory by \emph{byte}, not by word: consecutive
addresses differ by one byte. Words are built from bytes at aligned
addresses. So while a register moves one word at a time, an address
still names one byte---and pointer arithmetic from Chapter~1 works in
units of the pointed-to type, not in words. Both facts are true at
once: bytes are addressable, words are natural.
\end{insight}

\subsection{Registers, Cache, and RAM: The Memory Hierarchy}

Between the CPU and the huge, slow RAM we have been assuming sits
something your code never sees but that decides how fast it runs:
the \emph{memory hierarchy}. Access to a byte is not uniform. The CPU
keeps a handful of registers (fastest, tiny), then one or more
\emph{caches} (fast, medium-sized), then RAM (slower, huge), then
disk (far slower, enormous). Each level holds a copy of data that the
next level up is using, in blocks.

\begin{verbatim}
   The hierarchy, with the "distance from the CPU":

     level      size      speed        managed by
     -------    -------   ----------   ---------------
     registers   ~100s B   ~1 ns        the program (you!)
     L1 cache     ~32 KB    ~1 ns       the CPU, automatically
     L2 cache    ~256 KB    ~3 ns       the CPU, automatically
     L3 cache    ~8 MB      ~10 ns      the CPU, automatically
     RAM         ~GB        ~100 ns     the OS (the page model)
     disk        ~TB        ~ms         the OS (filesystems)

   (Numbers are indicative; what matters is the SHAPE: each level is
    bigger and slower, and faster levels cache the slower ones.)
\end{verbatim}

A \emph{cache} is a small, fast store that keeps recently used data so
the CPU rarely has to go to the slower level. When the CPU wants a
word, it looks in the cache first; only a \emph{cache miss} sends it
to RAM. Data moves between cache and RAM in fixed-size units called
\emph{cache lines}---typically 64 bytes. This one fact explains more
systems performance than any other:

\begin{verbatim}
   Why cache lines matter to your C code:

     two adjacent ints:  cache line = [ 42 ][ 43 ]...  one fetch loads
     both. A struct's fields arrive together (ch.3's AoS).

     two ints 1000 bytes apart: two cache lines, two fetches. If an
     algorithm walks field-by-field across many structs, it drags
     whole cache lines it never uses.

   Consequence: the COST of an operation depends on which cache lines
   it touches, not just how many instructions it runs. This is why
   Chapter 3's "array of structs vs struct of arrays" decision, and
   Chapter 4's unrolled lists, are performance choices, not style.
\end{verbatim}

The hierarchy also explains the numbers this book quotes. ``A cache
miss costs about 100 times a cache hit'' is the real reason aligned,
sequential, localized data wins. And the page model of the OS sits one
level down: RAM caches disk, exactly as L1 caches RAM. Same shape,
different scale.

\begin{insight}[Cache is invisible and irresistible]
The CPU manages the caches with no help from your program. You cannot
store into L1 ``by hand.'' But you \emph{can} arrange your data so
that the automatic caching works for you---contiguous, aligned, walked
in order. That arrangement is most of what ``performance engineering''
means, and it runs through every later chapter.
\end{insight}

\section{Sizes of the Primitive Types}

C's primitive types are contracts about the interpretation of a range of
bytes. The type tells the compiler (a) how many bytes the object
occupies, and (b) how to interpret those bytes. The standard defines
\emph{minimum} sizes, but on our target machines---a modern 64--bit
x86-64 or ARM64 system---the sizes are stable:

\begin{verbatim}
   type            size (bytes)   typical range
   --------------  ------------   ------------------------------------
   char            1              -128 .. 127           (signed)
   unsigned char   1                0 .. 255
   short           2             -32768 .. 32767
   unsigned short  2                0 .. 65535
   int             4              -2147483648 .. 2147483647
   unsigned int    4                0 .. 4294967295
   long            8              (on 64-bit Linux: 8 bytes)
   unsigned long   8
   float           4              ~7 decimal digits of precision
   double          8              ~15 decimal digits of precision
   pointer         8              (8 bytes, an address in the space)
   --------------  ------------   ------------------------------------
\end{verbatim}

Do not trust your memory of this table on every platform: in C, you
always ask the compiler for the true size at build time. That is the
job of \code{sizeof}.

\begin{lstlisting}[caption={Sizes as reported by the compiler itself.}]
#include <stdio.h>

int main(void) {
    printf("char                = %zu bytes\n", sizeof(char));
    printf("short               = %zu bytes\n", sizeof(short));
    printf("int                 = %zu bytes\n", sizeof(int));
    printf("long                = %zu bytes\n", sizeof(long));
    printf("double              = %zu bytes\n", sizeof(double));
    printf("int pointer         = %zu bytes\n", sizeof(int *));
    printf("void pointer        = %zu bytes\n", sizeof(void *));
    return 0;
}
\end{lstlisting}

Notice the format specifier \code{\%zu}: the letter \code{z} says
``the argument is a \code{size\_t}'', which is exactly the unsigned
type \code{sizeof} yields. You will meet \code{size\_t} constantly;
it is the type of every count in the standard library.

\begin{ruleofthumb}[The Loop Counter Rule]
In this course we use \code{size\_t} \emph{only} for loop counters that
iterate over arrays and buffers. Everywhere else---indices computed from
user input, signed arithmetic, difference between addresses---we use
plain \code{int} or \code{long}. This keeps signs explicit and avoids
the silent traps of unsigned wraparound when you subtract one index from
another.
\end{ruleofthumb}

\section{The Address Space of a Process}

When the operating system launches your program, it does not hand your
process raw physical RAM. It constructs a \emph{virtual address space}:
a private, contiguous-looking map of addresses in which your program
lives. Each process gets the illusion of owning the whole machine. The
map is laid out in a fixed order, from low addresses to high addresses:

\begin{verbatim}
   +--------------------------------------------------+  high addresses
   |                                                  |  0x7fffffffffff
   |                     STACK                        |
   |  local variables, return addresses, arguments    |
   |  grows DOWN toward lower addresses               |
   |  |                                               |
   |  v                                               |
   +--------------------------------------------------+
   |              (free region; this is where          |
   |               memory mappings and shared libs    |
   |               get placed)                        |
   +--------------------------------------------------+
   |  ^                                               |
   |  |                                               |
   |                     HEAP                         |
   |  memory handed out by malloc()/calloc()/free()  |
   |  grows UP toward higher addresses                |
   +--------------------------------------------------+
   |                     BSS                          |
   |  uninitialized globals and statics (zeroed)     |
   +--------------------------------------------------+
   |                     DATA                         |
   |  initialized globals, statics, string literals  |
   +--------------------------------------------------+
   |                     TEXT                         |
   |  the machine code of your program                |
   +--------------------------------------------------+  low addresses
   |                     0x0  (unmapped, fault)       |
   +--------------------------------------------------+
\end{verbatim}

This diagram is worth memorizing; we will return to it in every later
chapter. Three regions dominate our mental model:

\begin{description}
  \item[TEXT] Read-only. The instructions themselves. If you try to write
    here, the CPU raises a fault.
  \item[DATA/BSS] Globals and statics. BSS is not stored in the binary;
    the OS simply promises to give you zeroed memory at startup.
  \item[STACK] Automatic (local) variables, function call frames, return
    addresses. It grows \emph{down}.
  \item[HEAP] Everything you allocate at runtime with \code{malloc}.
    It grows \emph{up}.
\end{description}

The stack and the heap grow toward each other. This is why the phrase
``stack overflow'' exists: a runaway recursive function walks the stack
down until it collides with something it may not touch, and the
operating system kills the process. Not a metaphor; a literal collision
of addresses. (A \emph{recursive} function is simply one that calls
itself; each call pushes another frame onto the stack, and one that
never stops pushing is what ``overflows'' it. We write real recursive
functions---tree traversal, directory walking---starting in
Chapter~4.)

\subsection{Virtual Memory and Pages}

Two more words from the diagram need exact definitions now, because
they will carry every chapter from here on.

\emph{Virtual memory} is the mechanism behind the ``virtual address
space'' just described. Your program's addresses---the \code{0x7ffc}
numbers you printed in Chapter~1---are \emph{virtual}: they are not
physical RAM addresses at all. The kernel keeps a translation table,
the \emph{page table}, and the CPU consults it on every access,
converting your virtual address into the physical address where the
byte actually lives. This gives every process a private, huge address
space, \emph{isolation} (your page table never contains another
process's physical pages), and \emph{flexibility} (a virtual region can
be backed by RAM, by a file, or by nothing yet).

A \emph{page} is the unit of that machinery. Virtual memory is not
managed byte by byte; it is chopped into fixed-size chunks---4~KB on
our machines---and each chunk is mapped, protected, and shared as a
whole.

\begin{verbatim}
   Virtual memory, drawn:

   YOUR address space (virtual)        PHYSICAL RAM        the page table
   +-----------------+                +-----------+        maps virtual
   | stack (pages)   |                |           |        page -> physical
   | heap  (pages)   |  ---page---->  |           |        frame
   | data  (pages)   |  table lookup  |           |
   | text  (pages)   |                |           |
   +-----------------+                +-----------+

   a virtual address: 0x7ffc1004
        |
        | split by the CPU into:  page number  0x7ffc   +  offset 0x004
        |                         (look up in the page table)
        v
   physical address: the RAM cell that holds the byte

   - A page that is "not present" in the page table is not in RAM;
     touching it makes the CPU raise a PAGE FAULT, and the kernel
     decides what that means (load the page, or kill the process).
   - This is why the address space can look far larger than RAM: most
     pages are simply not present, and RAM holds only what is in use.
   - Each process has its OWN page table, which is what makes two
     processes using the same virtual address able to hold different
     data -- and what makes process A unable to read process B.
\end{verbatim}

The page explains several things you have already seen. A 4096-byte
buffer in Chapter~5 is exactly one page. The stack that ``overflows''
is a set of pages the kernel maps on demand and revokes on fault. And
the mmap machinery of Chapter~12 and the allocators of Chapter~13 are
both, at bottom, negotiations over pages.

\begin{verbatim}
   One page, drawn at scale (4 KB = 4096 bytes = 512 eight-byte words):

     +--------------------------------+  address P (page-aligned)
     |                                |
     |         4096 bytes             |
     |         (this page)            |
     |                                |
     +--------------------------------+  address P + 4095

   Pages are PAGE-ALIGNED: a page's address is a multiple of 4096.
   That is why malloc() returns 16-byte-aligned pointers but the heap
   itself is built from page-sized regions.
\end{verbatim}

\begin{insight}[A page is the OS's atom]
The OS never manages memory in bytes; it manages in pages. Everything
that costs you performance---a page fault, a copy-on-write, a
shared-memory mapping, a demand-paged file---happens page by page, at
4~KB granularity, and nothing you do in user space is ever smaller
than that at the OS level. When a program feels "slow on first
access," the page is the culprit.
\end{insight}

\section{Addresses and Pointers}

A \emph{pointer} is a variable whose value is an address. The type of a
pointer records two things: the kind of data stored at the destination,
and the fact that the value is an address. The type is a contract; the
value is just an integer that happens to be a number of a memory cell.

\begin{lstlisting}[caption={A pointer is a variable holding an address.}]
#include <stdio.h>

int main(void) {
    int value = 42;
    int *ptr;                 /* "ptr will hold the address of an int" */

    ptr = &value;             /* &value: take the address of value */

    printf("value  = %d\n", value);
    printf("&value = %p\n", (void *) &value);
    printf("ptr    = %p\n", (void *) ptr);
    printf("*ptr   = %d\n", *ptr);

    return 0;
}
\end{lstlisting}

The \code{\&} operator, called \emph{address-of}, produces the address
of the variable that follows it. The \code{*} operator, called
\emph{dereference}, goes from the pointer to the object it points at.
They are inverses: \code{*(\&value)} is \code{value}.

Let us draw the state of memory at the moment just before the final
\code{printf}:

\begin{verbatim}
   Suppose the compiler placed:
     value at address 0x7ffc1004
     ptr   at address 0x7ffc1008

   We depict each memory cell as a box containing its value.
   A pointer is drawn as an arrow to the cell it targets.

              address       bytes               interpretation
   value      0x7ffc1004    [ 42 ]              an int holding 42
   ptr        0x7ffc1008    [ 0x7ffc1004 ]      an int* holding value's address

   Box-and-arrow view:

        +------------------+          +------------------+
        |       ptr        |          |      value       |
        |  0x7ffc1008      |--------->|      42          |
        |  (holds address) |          |  (an int)        |
        +------------------+          +------------------+

   And in full:

        +---------------------------------------------+
        |  ptr = &value                               |
        +---------------------------------------------+
        |  "ptr"                |  "value"            |
        |  addr 0x7ffc1008      |  addr 0x7ffc1004    |
        |  value 0x7ffc1004 --->|  value 42           |
        +---------------------------------------------+
\end{verbatim}

The arrow is the single most important diagram in this book. When you
read C code, your job is to draw the arrows. When you write C code,
your job is to know exactly which arrows exist and where they point.
A C program is only correct if its arrows and its ownership rules are
correct; everything else follows.

\section{Dereferencing Is Reading or Writing Through an Arrow}

Dereferencing is a two-way street. It can \emph{read} the cell the arrow
points to, or it can \emph{write} through it. A dereference on the left
side of an assignment writes to the pointed-to cell.

\begin{lstlisting}[caption={Reading and writing through a pointer.}]
#include <stdio.h>

int main(void) {
    int value = 42;
    int *ptr = &value;

    /* Read: fetch the int that ptr points at. */
    printf("before: %d\n", *ptr);

    /* Write: store a new int INTO the cell ptr points at. */
    *ptr = 100;

    printf("after:  value = %d, *ptr = %d\n", value, *ptr);

    /* value was modified, because ptr points at value. */
    return 0;
}
\end{lstlisting}

\begin{verbatim}
   Before:                      After *ptr = 100;

     ptr ---> [ 42 ]              ptr ---> [ 100 ]
                 value                       value
\end{verbatim}

There is nothing magic here. \code{*ptr = 100} means: ``go to the
address stored in \code{ptr}, and write the bytes of \code{100} there.''
The C compiler turns it into one store instruction. The pointer is just
the messenger; the object it points at is what actually changes.

\begin{insight}[Pointers are how C returns multiple values]
Functions return one value. If you want a function to deliver two
results, you pass pointers to the caller's variables and the function
writes through them. This is exactly how \code{scanf} works, and how
we will write our own data structure APIs in later chapters:
\code{int list\_pop(struct list *l, int *out\_value);}
\end{insight}

\section{Pointer Arithmetic}

A pointer plus an integer steps over whole \emph{elements}, not bytes.
If \code{int *p} holds an address, then \code{p + 1} holds the address
of the next \code{int} object, which is 4 bytes later on our machine.
The compiler knows the size of the pointed-to type, so it does the
multiplication for you.

This is the C idiom that builds every array algorithm, and it is the
place where ``clever'' code tends to hide. In this course we refuse
obfuscation: we move pointers on explicit lines, we keep the current
element in a named variable, and we only ever use \code{size\_t} as the
loop counter. Compare the same loop written two ways.

The style you will find in legacy code (and, alas, in many textbooks):

\begin{lstlisting}[caption={Legacy style -- compact, but opaque.}]
/* Copies the string src into the buffer dst. Obscured, hard to debug. */
void copy_string(char *dst, const char *src) {
    while (*dst++ = *src++);   /* we do NOT teach this */
}
\end{lstlisting}

That line does four operations inside a test. If a bug appears, you
cannot step through it one move at a time. The style we use:

\begin{lstlisting}[caption={Our explicit style: one move per line, named pointers.}]
#include <stddef.h>
#include <string.h>

/* Copy src into dst, including the terminating NUL byte. */
/* Precondition: dst has room for strlen(src) + 1 bytes.     */
void copy_string(char *dst, const char *src) {
    size_t i;
    size_t len;

    len = strlen(src);

    for (i = 0; i <= len; i++) {
        /* Copy one byte: the element at position i of src. */
        dst[i] = src[i];
    }
}
\end{lstlisting}

Notice what we gained: every step is a single, auditable statement.
We can put a breakpoint on any line, print \code{i} and \code{len},
and inspect the bytes one at a time. The clarity costs almost nothing
in speed; a modern compiler turns the explicit loop into the same
optimized code it would generate for the clever version. Clarity is
nearly free. We take it.

\section{Arrays Are Contiguous Bytes; Pointers Point at Elements}

An array of \code{N} elements of type \code{T} is simply \code{N}
objects of type \code{T} laid out consecutively in memory. There is no
header, no length, no metadata. The address of the array is the address
of its first element.

\begin{verbatim}
   int numbers[4] = {10, 20, 30, 40};

   A 4-byte int is drawn as one box. The array occupies 16 consecutive
   bytes:

      numbers +----+----+----+----+
              | 10 | 20 | 30 | 40 |
              +----+----+----+----+
      address: 1000 1004 1008 1012

      numbers            == &numbers[0]   == 1000
      numbers + 1        == &numbers[1]   == 1004
      numbers + 2        == &numbers[2]   == 1008

      *numbers           == 10
      *(numbers + 2)     == 30            == numbers[2]
\end{verbatim}

The expression \code{numbers[i]} is defined to mean
\code{*(numbers + i)}. The name of an array, used in an expression,
\emph{decays} into a pointer to its first element. This is why
passing an array to a function passes the address, not a copy: any
write inside the function modifies the caller's array.

\begin{lstlisting}[caption={Arrays decay to pointers when passed to functions.}]
#include <stdio.h>

/* Fill a buffer of n ints with the values 0..n-1. */
void fill_seq(int *buf, size_t n) {
    size_t i;

    for (i = 0; i < n; i++) {
        buf[i] = (int) i;
    }
}

int main(void) {
    int numbers[4] = {0, 0, 0, 0};
    size_t i;

    fill_seq(numbers, 4);           /* "numbers" is &numbers[0] here */

    for (i = 0; i < 4; i++) {
        printf("numbers[%zu] = %d\n", i, numbers[i]);
    }

    return 0;
}
\end{lstlisting}

Note the parameter type \code{int *buf}. In the declaration of a
function parameter, \code{int buf[]} and \code{int *buf} are identical.
We write the pointer form, because it tells the truth about what is
passed.

\section{Stack vs.\ Heap: Who Owns the Memory?}

This distinction is the pivot of all of C. Two places hold the data
your program manipulates, and they differ in \emph{lifetime} and
\emph{management}.

\subsection{The Stack}

Local variables live on the stack. When a function is called, the
machine pushes a \emph{stack frame} containing its locals and return
address; when the function returns, that frame is popped and its
storage is, effectively, gone. Stack memory is:

\begin{itemize}
  \item allocated and freed automatically (by the call/return protocol),
  \item extremely fast,
  \item limited in size (typically ~8 MB on Linux),
  \item invalid the moment the function returns.
\end{itemize}

\begin{lstlisting}[caption={A classic lifetime bug: returning a pointer to a stack local.}]
#include <stdio.h>

int *bad_function(void) {
    int local = 42;

    return &local;   /* ERROR: local dies when bad_function returns */
}

int main(void) {
    int *p = bad_function();

    /* p dangles: it names memory that the stack has already reclaimed. */
    printf("%d\n", *p);   /* undefined behavior */
    return 0;
}
\end{lstlisting}

\begin{verbatim}
   The stack during main() -> bad_function() -> return:

      before call:     during bad_function():    after return:
      +--------------+  +--------------+          +--------------+
      | main frame   |  | main frame   |          | main frame   |
      |  p = ?       |  |  p = ?       |   p ---> |  p = 0x...   |
      +--------------+  +--------------+          +--------------+
      (stack top)       | bad frame    |          (frame popped)
                        |  local = 42  |<-- &local is now garbage
                        +--------------+
\end{verbatim}

The address value survives; the storage behind it does not. We call
\code{p} a \emph{dangling pointer}. No compiler can reliably catch
every use of one; only discipline can. The rule:

\begin{ruleofthumb}[Heap for anything that outlives its creator]
If data must be returned from a function, or shared between functions,
or sized by values known only at runtime, it goes on the \emph{heap},
allocated with \code{malloc}, and freed with \code{free} by whichever
party becomes its owner.
\end{ruleofthumb}

\subsection{The Heap}

The heap is managed explicitly. \code{malloc(size)} asks the OS (via a
bigger allocator such as \code{glibc}) for \code{size} bytes and
returns a pointer to them, or \code{NULL} if it cannot. The bytes are
\emph{yours} until you call \code{free}. Nothing is reclaimed
automatically. This is the price of lifetime control, and it is a fair
price: with it we build linked lists, hash tables, trees, and whole
servers, all with deterministic behavior.

\begin{lstlisting}[caption={A correct heap round-trip.}]
#include <stdio.h>
#include <stdlib.h>
#include <string.h>

int main(void) {
    char *buffer;
    size_t len;

    len = 32;
    buffer = malloc(len);
    if (buffer == NULL) {
        fprintf(stderr, "out of memory\n");
        exit(1);
    }

    strcpy(buffer, "hello from the heap");

    printf("%s\n", buffer);

    free(buffer);          /* hand the bytes back */

    /* After free, buffer is a dangling pointer; do not dereference it. */
    return 0;
}
\end{lstlisting}

The \code{== NULL} check is not optional. \code{malloc} fails when the
system cannot honor the request, and a good program treats that
possibility as part of its normal control flow.

\section{Pointers to Pointers}

Sometimes the thing a pointer points at is, itself, a pointer. A
\code{char **argv} is a pointer to an array of strings; a function that
must \emph{reassign} a caller's pointer takes a pointer to that pointer.
This is how we implement functions that ``return'' newly allocated
objects through their arguments.

\begin{lstlisting}[caption={Writing a new pointer into a caller's variable.}]
#include <stdio.h>
#include <stdlib.h>

/* Allocates a small string and stores its address in *out. */
void make_message(char **out) {
    char *s = malloc(16);

    if (s == NULL) {
        *out = NULL;
        return;
    }
    s[0] = 'O'; s[1] = 'K'; s[2] = '\0';

    *out = s;          /* write the pointer into the caller's variable */
}

int main(void) {
    char *msg = NULL;

    make_message(&msg);      /* &msg is a char** */

    if (msg != NULL) {
        printf("%s\n", msg);
        free(msg);
    }

    return 0;
}
\end{lstlisting}

\begin{verbatim}
   Before make_message(&msg):        After *out = s:

     msg  [ NULL ]                     msg  [ 0x.. ] ---> [ O ][ K ][ 0 ]
     out  [ &msg ]                     out  [ &msg ]
\end{verbatim}

The pattern is always the same: the parameter is a pointer to the
variable we want the function to modify, and inside the function we
dereference it to perform the write.

\section{NULL: The Pointer That Points at Nothing}

\code{NULL} is the address 0, and by convention it means ``this pointer
is not currently valid''. It is the only address you may
\emph{compare} against but never \emph{dereference}. The processor does
not map address 0, so dereferencing a \code{NULL} pointer traps
immediately with a segmentation fault, which is a favor: the bug
surfaces exactly where it happens, instead of corrupting data silently
a million instructions later.

\begin{verbatim}
   ptr ---> [ 42 ]       valid pointer, points at a real int

   ptr = NULL;

   ptr ---> 0x0          NOT a real object.
   Dereferencing it:    *ptr  ==>  SIGSEGV (segmentation fault)
\end{verbatim}

Use \code{NULL} honestly: every allocation is checked against it,
every function that may fail returns it, and every pointer in a
structure that is ``not set'' holds it. \code{NULL} is the C equivalent
of ``no object here'', and it costs nothing to check but a comparison.

\section{Undefined Behavior: The Abyss}

The C standard reserves the right to do \emph{anything} when your
program crosses certain lines: dereferencing a dangling or \code{NULL}
pointer, reading an uninitialized variable, overflowing a signed
integer, indexing an array out of bounds. The compiler may assume you
never do these things, and optimize accordingly. When you violate the
contract, the program does not ``throw an exception''; it simply has no
guaranteed behavior. It might work today and break tomorrow, or under a
different compiler flag, or when the array happens to sit next to
something important.

This is why the discipline in this course is non-negotiable:

\begin{itemize}
  \item Every pointer is either valid, \code{NULL}, or untouched.
  \item Every array index is checked against a length you carry.
  \item Every \code{malloc} result is checked.
  \item Every byte a function writes is justified by a size argument.
\end{itemize}

\section{Command-Line Arguments: argc and argv}

Almost every real program is launched with arguments: \code{./tool
input.txt -v}. Those arguments arrive in \code{main} through two
parameters, and since every later chapter's practical programs use
them, we define them once here.

\begin{lstlisting}[caption={main receives its command line as argc and argv.}]
#include <stdio.h>

int main(int argc, char *argv[]) {
    int i;

    printf("program name: %s\n", argv[0]);
    printf("number of arguments: %d\n", argc);

    for (i = 1; i < argc; i++) {
        printf("  argv[%d] = %s\n", i, argv[i]);
    }

    return 0;
}
\end{lstlisting}

\begin{verbatim}
   Running:   ./prog alpha "two words"

   argc = 3           (the count of strings in argv)
   argv[0] = "./prog"       the program's own name
   argv[1] = "alpha"
   argv[2] = "two words"    (quotes group a string with a space)

   main's parameters:
     argc   an int: how many strings are in argv (always >= 1)
     argv   an array of char* strings; argv[argc] is always NULL

   - The SHELL (or your parent process) builds this array and passes
     it to the new program at exec time (ch.5). The kernel copies it
     onto the new program's stack.
   - argv[0] is not an argument the user typed; it is the name the
     program was invoked as. Arguments the user typed begin at argv[1].
   - That is why every practical in this book guards with
     "if (argc < 2) { fprintf(stderr, \"usage...\"); }": it is
     checking that the user actually supplied argv[1].
\end{verbatim}

The two spellings \code{char *argv[]} and \code{char **argv} are
identical here, by the array-decay rule of Section~1.8. We write
\code{char *argv[]} because it says what it is: an array of strings.
And the strings themselves---the \code{argv[i]} values---are ordinary
\code{char*} pointers into a region the kernel placed on the stack
before \code{main} ran.

\section{Practical Exercise: A Memory Inspector}
Let us consolidate the chapter into a small, real tool. It prints the
layout of a few variables and then hex-dumps a region of memory so you
can literally \emph{see} bytes, addresses, and pointer values. Run it;
look at the output until it stops being surprising.

\begin{lstlisting}[caption={memory\_inspect.c: see bytes, addresses, and pointers.}]
#include <stdio.h>
#include <stdlib.h>

/* Prints the raw bytes of the n bytes starting at p. */
void hex_dump(const unsigned char *p, size_t n) {
    size_t i;

    for (i = 0; i < n; i++) {
        /* Print byte number i, then a separator between groups of 8. */
        printf("%02x", p[i]);
        if ((i % 8) == 7) {
            printf("  ");
        } else {
            printf(" ");
        }
    }
    printf("\n");
}

int main(void) {
    int value = 0x11223344;   /* one int, four bytes we can recognize */
    int *ptr = &value;
    char text[] = "abc";      /* four bytes: 'a','b','c','\0'         */

    printf("=== addresses ===\n");
    printf("&value     = %p\n", (void *) &value);
    printf("&ptr       = %p\n", (void *) &ptr);
    printf("&text      = %p\n", (void *) &text);
    printf("ptr        = %p\n", (void *) ptr);

    printf("\n=== sizes ===\n");
    printf("sizeof(int)     = %zu\n", sizeof(int));
    printf("sizeof(int *)   = %zu\n", sizeof(int *));
    printf("sizeof(text)    = %zu\n", sizeof(text));

    printf("\n=== the int as raw bytes ===\n");
    hex_dump((unsigned char *) &value, sizeof(value));

    printf("\n=== the pointer as raw bytes ===\n");
    hex_dump((unsigned char *) &ptr, sizeof(ptr));

    printf("\n=== the string (with its NUL terminator) ===\n");
    hex_dump((unsigned char *) text, sizeof(text));

    printf("\n=== stepping over an int with pointer arithmetic ===\n");
    {
        /* p moves from value to each byte-sized window inside it. */
        unsigned char *q = (unsigned char *) &value;
        size_t i;

        for (i = 0; i < sizeof(value); i++) {
            printf("byte %zu at %p = 0x%02x\n", i, (void *) &q[i], q[i]);
        }
    }

    return 0;
}
\end{lstlisting}

If your machine is little-endian (x86-64 and ARM64 both are), the int
\code{0x11223344} is stored as the byte sequence \code{44 33 22 11}:
the least significant byte comes first in memory. The pointer's raw
bytes encode the address, and you will recognize them as the value
\code{ptr} printed in the addresses section.

\begin{insight}[Endianness]
Endianness is the order in which the bytes of a multi-byte number are
stored. Little-endian machines store the least significant byte first.
Most of the time you will not care, but the moment you write a file
format, a network protocol, or a virtual machine, you will care a great
deal. We will revisit endianness when we serialize data to disk and to
the network in later chapters.
\end{insight}

\begin{researchnote}[Memory allocators and custom allocators]
All through this chapter we told you to trust \code{malloc} and check
its result. Researchers have actually measured whether programs benefit
from \emph{custom allocators}---hand-written replacements for
\code{malloc} tuned for one program. ``Recent research'' here means a
careful empirical study by Emery Berger and colleagues, and its
conclusion is an excellent corrective to a very common habit: writing a
custom allocator is usually wasted effort.

\begin{verbatim}
   The study:  Berger, Zorn, and McKinley studied eight applications that use
   custom allocators and compared each one, in wall-clock time and in
   memory, against the Lea allocator (the general-purpose malloc of the
   day). Their finding, in one line:

       "Custom allocators are rarely effective: for most programs a
        general-purpose malloc performs as well as or better than the
        custom allocator; only region-based allocators (arenas)
        consistently show gains, of up to 44%."

   (That is a brief paraphrase of the paper's findings, condensed for the book.)
\end{verbatim}

What makes the study trustworthy, and worth quoting: it is an
\emph{empirical} result (measured, not argued), it was published in a
peer-reviewed venue, and it runs against the folklore that every
performance-sensitive program needs a hand-tuned allocator. The lesson
for this book is simple: use \code{malloc}, use \code{free}, check
failure, and profile \emph{before} you ever consider writing a custom
allocator. We will revisit this when we design our own allocators in
Chapter~13, where the goal will be \emph{policy} (object pools, arenas)
rather than raw speed.

\paperref{E.~D.~Berger, B.~G.~Zorn, K.~S.~McKinley, ``Reconsidering
Custom Memory Allocation,'' ACM OOPSLA 2002}{https://people.cs.umass.edu/~emery/pubs/berger-oopsla2002.pdf}
\end{researchnote}

\section{Structures in Bytes}

Every structure in this book---variable, struct, packet, inode---is,
at bottom, a fixed arrangement of bytes. This section starts the habit
of reading that arrangement; every later chapter ends with its own
byte map. On our 64-bit little-endian machine, the scalars look like
this:

\begin{verbatim}
   type         size   bytes (little-endian storage)
   ---------    ----   ------------------------------------------------
   char         1      01
   short        2      01 00
   int          4      01 00 00 00
   long         8      01 00 00 00 00 00 00 00
   float        4      00 00 80 3f            (the float value 1.0)
   double       8      00 00 00 00 00 00 f0 3f   (the double value 1.0)
   pointer      8      e6 3c a2 29 a0 00 00 00   (some address)
\end{verbatim}

Little-endian means byte 0 is the least significant byte, so the value
\code{1} is stored as \code{01 00 00 00} and \code{0x11223344} as
\code{44 33 22 11}. Reading a memory dump is reading these sequences in
groups of the type's size. A stack frame is bytes too: on x86-64 the
\code{call} instruction pushes an 8-byte return address, the callee
saves the old frame pointer, then makes room for its locals:

\begin{verbatim}
   one stack frame (the stack grows downward):

   +-------------------------------+
   | caller's frame                |
   +-------------------------------+
   | return address          (8B)  |  pushed by the call instruction
   +-------------------------------+
   | saved frame pointer     (8B)  |  pushed by "push rbp"
   +-------------------------------+  <- rbp
   | locals (e.g. int value, 4B)   |
   | and saved registers           |
   +-------------------------------+  <- rsp
\end{verbatim}

\code{\&value} is a byte address inside that range---which is why the
frame diagram of Chapter~1 and the byte diagram here are the same
picture.

\section{Summary of the Mental Model}

At the end of this chapter you should be able to answer, for any
snippet of C you read: \emph{which variable lives where, how many bytes
it occupies, what pointer points at what, and who owns the memory.}

\begin{verbatim}
   The four questions of every C program:

   1. Address : where do the bytes live?       (stack / heap / globals)
   2. Size    : how many bytes?                (sizeof gives the truth)
   3. Arrows  : which pointers point where?    (draw the diagram)
   4. Owner   : who calls free()?              (exactly one party)
\end{verbatim}

If you can draw the memory diagram for a program, you understand the
program. If you cannot, you do not---regardless of how many tests pass.
This habit, more than any feature of the language, is what separates
system programmers from everyone else.

In the next chapter we open the machine's fastest memory: the
\emph{registers}, built from flip-flops---bistable loops of
transistors that hold a bit---and descend from the register image to
the electronics and chemistry that make it retain its value. The
chapter after that combines scalars into \emph{structures}, and
learns why the machine, in its demand for aligned access, quietly
inserts invisible bytes between our fields.
